% Options for packages loaded elsewhere
\PassOptionsToPackage{unicode}{hyperref}
\PassOptionsToPackage{hyphens}{url}
\documentclass[
]{article}
\usepackage{ma2003b-notes}
\hypersetup{
  pdftitle={Multivariate Regression - Lecture Notes},
  pdfauthor={MA2003B},
  hidelinks,
  pdfcreator={LaTeX via pandoc}}

\title{Multivariate Regression - Lecture Notes}
\author{MA2003B}
\date{}

\begin{document}
\maketitle

\textbf{Multivariate Regression}

Lecture Notes

MA2003B - Application of Multivariate Methods in Data Science

\begin{center}\rule{0.5\linewidth}{0.5pt}\end{center}

\textbf{ }

\section{Document Information}

\textbf{ }

\subsection{Credits and Roles}

\textbf{Instructional Architect: Juliho Castillo Colmenares (Course
Instructor)}

\begin{itemize}
\item
  Designed course structure and learning objectives
\item
  Established pedagogical direction and content scope
\item
  Defined assessment criteria and evaluation approach
\item
  Reviewed content for academic accuracy and rigor
\item
  Maintains ongoing responsibility for course materials
\end{itemize}

\textbf{AI Assistant: Claude Sonnet 4.5 (Anthropic)}

\begin{itemize}
\item
  Generated comprehensive lecture notes based on curriculum design
\item
  Developed mathematical formulations and technical explanations
\item
  Created practice questions and detailed answers
\item
  Structured content according to pedagogical best practices
\item
  Aligned material with E07 Multivariate Regression quiz topics
\end{itemize}

\textbf{ }

\subsection{Contact Information}

\textbf{Course Instructor:}

\begin{itemize}
\item
  \textbf{Name:} Juliho Castillo Colmenares
\item
  \textbf{Email:} julihocc@tec
\item
  \textbf{Office:} Tec de Monterrey CCM Office 1540
\item
  \textbf{Office Hours:} Monday-Friday, 9:00 AM - 5:00 PM
\end{itemize}

\textbf{Course Information:}

\begin{itemize}
\item
  \textbf{Course Code:} MA2003B
\item
  \textbf{Course Title:} Application of Multivariate Methods in Data
  Science
\item
  \textbf{Institution:} Tec de Monterrey
\item
  \textbf{Department:} Escuela de Ingeniería y Ciencias
\end{itemize}

\textbf{ }

\subsection{Version Information}

\begin{itemize}
\item
  \textbf{Document Version:} 1.0
\item
  \textbf{Created:} 2025-11-12
\item
  \textbf{Last Updated:} 2025-11-12
\item
  \textbf{Status:} Active - Current Semester
\end{itemize}

\textbf{ }

\subsection{Usage and Distribution}

These lecture notes are intended for students enrolled in MA2003B. The
material is designed to prepare students for the Multivariate Regression
evaluation (E07) and provide comprehensive understanding of advanced
regression methods in multivariate settings.

\textbf{Academic Integrity:} Students are expected to use these notes as
a study resource in accordance with institutional academic integrity
policies.

\begin{center}\rule{0.5\linewidth}{0.5pt}\end{center}

\begin{center}\rule{0.5\linewidth}{0.5pt}\end{center}

\textbf{ }

\section{Logistic Regression Model}

\textbf{ }

\subsection{Introduction to Logistic Regression}

Logistic regression is a statistical method for modeling binary response
variables. Unlike linear regression, which models continuous outcomes,
logistic regression predicts the probability that an observation belongs
to one of two categories.

\textbf{Info: } The response variable in logistic regression follows a
\textbf{Bernoulli distribution}, meaning it can take only two values: 0
(failure/absence) or 1 (success/presence).

\textbf{ }

\subsection{The Logit Link Function}

The core of logistic regression is the \textbf{logit function}, which
transforms probabilities to the real line:

\[\text{ logit}(p) = \log(\frac{p}{1 - p}) = \beta_{0} + \beta_{1}x_{1} + \ldots + \beta_{k}x_{k}\]

Where:

\begin{itemize}
\item
  \(p\) is the probability of the event occurring
\item
  \(\frac{p}{1 - p}\) is the \textbf{odds} of the event (a ratio
  compares two odds, e.g. \(e^{\beta_{j}}\) below is the odds
  \textbf{ratio} per unit increase in \(x_{j}\))
\item
  \(\beta_{0},\beta_{1},\ldots,\beta_{k}\) are regression coefficients
\end{itemize}

The inverse transformation gives us the predicted probability:

\[p = \frac{1}{1 + \exp( - \left( \beta_{0} + \beta_{1}x_{1} + \ldots + \beta_{k}x_{k} \right))}\]

\textbf{ }

\subsection{Key Differences from Linear Regression}

\textbf{Warning: } Do not use ordinary least squares (OLS) for binary
outcomes. The assumptions of linear regression are violated when the
response is binary:

\begin{itemize}
\item
  Residuals are not normally distributed
\item
  Variance is not constant (heteroscedasticity)
\item
  Predicted values can fall outside 0,1 range
\end{itemize}

Logistic regression addresses these issues by:

\begin{enumerate}
\item
  Using a link function to bound predictions between 0 and 1
\item
  Modeling probabilities rather than raw values
\item
  Using maximum likelihood estimation instead of least squares
\end{enumerate}

\textbf{ }

\subsection{Maximum Likelihood Estimation}

Since OLS assumptions are violated, we use \textbf{Maximum Likelihood
Estimation (MLE)} to estimate coefficients. The likelihood function for
binary data is:

\[L(\beta) = \prod_{i = 1}^{n}p_{i}^{y_{i}}\left( 1 - p_{i} \right)^{1 - y_{i}}\]

Where \(y_{i}\) is the observed outcome (0 or 1) and \(p_{i}\) is the
predicted probability.

The log-likelihood is maximized using iterative numerical methods
(typically Newton-Raphson or Fisher scoring).

\textbf{ }

\subsection{Interpretation of Coefficients}

In logistic regression, coefficients represent the change in log-odds
for a one-unit increase in the predictor:

\begin{itemize}
\item
  \(\beta_{j} > 0\): Increasing \(x_{j}\) increases the odds of the
  event
\item
  \(\beta_{j} < 0\): Increasing \(x_{j}\) decreases the odds of the
  event
\end{itemize}

\textbf{Tip: } To interpret effects on probability scale, compute the
odds ratio: \(\text{OR } = \exp(\beta_{j})\). An OR greater than 1
indicates increased odds, while OR less than 1 indicates decreased odds.

\textbf{ }

\subsection{Model Diagnostics}

Common diagnostics for logistic regression include:

\begin{itemize}
\item
  \textbf{Deviance}: Measures goodness of fit
\item
  \textbf{AIC/BIC}: Model selection criteria
\item
  \textbf{ROC Curve and AUC}: Classification performance
\item
  \textbf{Hosmer-Lemeshow test}: Calibration assessment
\end{itemize}

\textbf{Example: } Predicting customer churn:
\(p\left( \text{churn} \right) = \frac{1}{1 + \exp( - \left( 2.5 - 0.8 \times \text{ tenure } + 1.2 \times \text{ complaints} \right))}\)

\begin{itemize}
\item
  Coefficient for tenure: -0.8 (longer tenure reduces churn probability)
\item
  Coefficient for complaints: 1.2 (more complaints increase churn
  probability)
\end{itemize}

\begin{center}\rule{0.5\linewidth}{0.5pt}\end{center}

\textbf{ }

\section{Inferences for Variances and Covariances Matrices}

\textbf{ }

\subsection{Multivariate Normal Distribution}

Before discussing inference for covariance matrices, we must establish
that the data comes from a multivariate normal distribution:

\[\mathbf{X} \sim N_{p}\left( \mathbf{\mu},\mathbf{\Sigma} \right)\]

Where:

\begin{itemize}
\item
  \(\mathbf{X}\) is a \(p \times 1\) random vector
\item
  \(\mathbf{\mu}\) is the mean vector
\item
  \(\mathbf{\Sigma}\) is the \(p \times p\) covariance matrix
\end{itemize}

\textbf{ }

\subsection{Sample Covariance Matrix}

For a sample of \(n\) observations, the unbiased estimate of the
covariance matrix is:

\[\mathbf{S} = \frac{1}{n - 1}\sum_{i = 1}^{n}\left( \mathbf{X}_{i} - \overline{\mathbf{X}} \right)\left( \mathbf{X}_{i} - \overline{\mathbf{X}} \right)'\]

\textbf{ }

\subsection{The Wishart Distribution}

The Wishart distribution is the multivariate generalization of the
chi-square distribution. When \(\mathbf{X}_{1},\ldots,\mathbf{X}_{n}\)
are independent \(N_{p}\left( \mathbf{\mu},\mathbf{\Sigma} \right)\):

\[(n - 1)\mathbf{S} \sim W_{p}\left( n - 1,\mathbf{\Sigma} \right)\]

This is read as ``\((n - 1)S\) follows a Wishart distribution with
\(n - 1\) degrees of freedom and scale matrix \(\Sigma\).''

\textbf{Info: } Just as the chi-square distribution is used for
inference about a single variance, the Wishart distribution is used for
inference about covariance matrices.

\textbf{ }

\subsection{Testing Equality of Covariance Matrices}

A common question is whether two or more groups have the same covariance
structure.

\subsubsection{Box's M Test}

For testing
\(H_{0}:\mathbf{\Sigma}_{1} = \mathbf{\Sigma}_{2} = \ldots = \mathbf{\Sigma}_{g}\),
Box's M statistic is:

\[M = (n - g)\log|\mathbf{S}_{\text{pooled}}| - \sum_{i = 1}^{g}\left( n_{i} - 1 \right)\log|\mathbf{S}_{i}|\]

Where:

\begin{itemize}
\item
  \(\mathbf{S}_{\text{pooled}}\) is the pooled covariance matrix
\item
  \(\mathbf{S}_{i}\) is the covariance matrix for group \(i\)
\item
  \(n_{i}\) is the sample size for group \(i\)
\end{itemize}

The raw statistic \(M\) is \textbf{not} itself chi-square distributed
and has no universal cutoff (e.g. ``\(M < 30\)'' is not a valid decision
rule). It must first be rescaled by a correction factor \(c_{1}\) (Box,
1949):

\[c_{1} = \left( \sum_{i = 1}^{g}\frac{1}{n_{i} - 1} - \frac{1}{n - g} \right) \cdot \frac{2p^{2} + 3p - 1}{6(p + 1)(g - 1)}\]

\[\chi^{2} = M\left( 1 - c_{1} \right),\quad\text{ df } = (g - 1)\frac{p(p + 1)}{2}\]

Only this corrected \(\chi^{2}\) statistic is compared against the
chi-square distribution to obtain a p-value. A non-significant result
means the test \textbf{fails to reject} equal covariance matrices -- it
does not prove equality.

\textbf{Warning: } Box's M test is sensitive to departures from
multivariate normality. Use with caution and verify normality
assumptions first.

\subsubsection{Bartlett's Test for Univariate Variances}

For the special case of testing equality of variances in univariate
data:

\[H_{0}:\sigma_{1}^{2} = \sigma_{2}^{2} = \ldots = \sigma_{g}^{2}\]

Bartlett's test statistic is:

\[\chi^{2} = \frac{(n - g)\log(s_{\text{pooled}}^{2}) - \sum_{i = 1}^{g}\left( n_{i} - 1 \right)\log(s_{i}^{2})}{C}\]

Where \(C\) is a correction factor for small samples.

\textbf{ }

\subsection{Assumptions for Valid Inference}

For valid inference about covariance matrices:

\begin{enumerate}
\item
  Data must come from a multivariate normal distribution
\item
  Observations must be independent
\item
  Sample sizes should be adequate (general rule: \(n > 5p\))
\end{enumerate}

\textbf{Tip: } Always test for multivariate normality before conducting
inference on covariance matrices. Methods include Mardia's test or
examining Q-Q plots for each variable.

\begin{center}\rule{0.5\linewidth}{0.5pt}\end{center}

\textbf{ }

\section{Inferences for a Vector of Means}

\textbf{ }

\subsection{Hotelling's T-Squared Test}

Hotelling's \(T^{2}\) statistic is the multivariate generalization of
the univariate t-test. It tests hypotheses about mean vectors.

\subsubsection{One-Sample Test}

For testing \(H_{0}:\mathbf{\mu} = \mathbf{\mu}_{0}\) against
\(H_{1}:\mathbf{\mu} \neq \mathbf{\mu}_{0}\):

\[T^{2} = n\left( \overline{\mathbf{X}} - \mathbf{\mu}_{0} \right)'\mathbf{S}^{- 1}\left( \overline{\mathbf{X}} - \mathbf{\mu}_{0} \right)\]

Under \(H_{0}\), the test statistic can be transformed to an
F-distribution:

\[F = \frac{n - p}{(n - 1)p}T^{2} \sim F_{p,n - p}\]

Where:

\begin{itemize}
\item
  \(n\) is sample size
\item
  \(p\) is number of variables
\item
  \(\overline{\mathbf{X}}\) is the sample mean vector
\item
  \(\mathbf{S}\) is the sample covariance matrix
\end{itemize}

\subsubsection{Two-Sample Test}

For comparing two independent groups, testing
\(H_{0}:\mathbf{\mu}_{1} = \mathbf{\mu}_{2}\):

\[T^{2} = \frac{n_{1}n_{2}}{n_{1} + n_{2}}\left( {\overline{\mathbf{X}}}_{1} - {\overline{\mathbf{X}}}_{2} \right)'\mathbf{S}_{\text{pooled}}^{- 1}\left( {\overline{\mathbf{X}}}_{1} - {\overline{\mathbf{X}}}_{2} \right)\]

The pooled covariance matrix is:

\[\mathbf{S}_{\text{pooled }} = \frac{\left( n_{1} - 1 \right)\mathbf{S}_{1} + \left( n_{2} - 1 \right)\mathbf{S}_{2}}{n_{1} + n_{2} - 2}\]

The F-transformation is:

\[F = \frac{n_{1} + n_{2} - p - 1}{\left( n_{1} + n_{2} - 2 \right)p}T^{2} \sim F_{p,n_{1} + n_{2} - p - 1}\]

\textbf{Info: } The key advantage of Hotelling's \(T^{2}\) over multiple
univariate t-tests is that it controls the overall Type I error rate
when testing multiple variables simultaneously.

\textbf{ }

\subsection{Paired Samples}

For paired observations, test \(H_{0}:\mathbf{\mu}_{D} = \mathbf{0}\)
where \(\mathbf{D} = \mathbf{X}_{1} - \mathbf{X}_{2}\):

\[T^{2} = n\overline{\mathbf{D}}'\mathbf{S}_{D}^{- 1}\overline{\mathbf{D}}\]

This reduces to the one-sample case applied to difference vectors.

\textbf{ }

\subsection{Confidence Regions}

Simultaneous confidence intervals for mean vectors form ellipsoidal
regions in \(p\)-dimensional space:

\[\left( \mathbf{\mu} - \overline{\mathbf{X}} \right)'\left( \frac{\mathbf{S}}{n} \right)^{- 1}\left( \mathbf{\mu} - \overline{\mathbf{X}} \right) \leq \frac{p(n - 1)}{n - p}F_{p,n - p,\alpha}\]

\textbf{Example: } Comparing nutritional content (protein, fat,
carbohydrates) between two diets using two-sample Hotelling's \(T^{2}\)
test. This accounts for correlations between nutritional components and
provides a single omnibus test.

\textbf{ }

\subsection{Rejection Criteria}

Reject \(H_{0}\) when:

\begin{itemize}
\item
  \(T^{2}\) exceeds the critical value from the \(T^{2}\) distribution,
  or equivalently
\item
  \(F\) exceeds the critical value \(F_{p,n - p,\alpha}\)
\end{itemize}

\textbf{Tip: } After rejecting the null hypothesis with Hotelling's
\(T^{2}\), follow up with univariate tests or discriminant analysis to
identify which variables contribute most to the difference.

\begin{center}\rule{0.5\linewidth}{0.5pt}\end{center}

\textbf{ }

\section{Multivariate Analysis of Variance (MANOVA)}

\textbf{ }

\subsection{Introduction to MANOVA}

\textbf{MANOVA} (Multivariate Analysis of Variance) extends ANOVA to
situations with multiple dependent variables. It tests whether group
means on a combination of dependent variables differ across levels of
categorical independent variables.

\textbf{ }

\subsection{Why Use MANOVA Instead of Multiple ANOVAs?}

\textbf{Warning: } Conducting separate ANOVAs for each dependent
variable increases Type I error through multiple testing. If you test 5
variables at alpha equals 0.05, your overall error rate could be as high
as 1 minus 0.95 raised to 5 equals 0.226.

Advantages of MANOVA:

\begin{enumerate}
\item
  \textbf{Controls Type I error} by testing all variables simultaneously
\item
  \textbf{Accounts for correlations} between dependent variables
\item
  \textbf{Can increase statistical power} by combining information
  across correlated variables -- the gain depends on the effect pattern,
  covariance structure, sample size, and multiplicity strategy; it is
  not guaranteed
\item
  \textbf{Detects patterns} that univariate tests might miss
\end{enumerate}

\textbf{ }

\subsection{MANOVA Model}

The multivariate linear model:

\[\mathbf{Y}_{n \times p} = \mathbf{X}_{n \times k}\mathbf{B}_{k \times p} + \mathbf{E}_{n \times p}\]

Where:

\begin{itemize}
\item
  \(\mathbf{Y}\) is the matrix of responses (\(n\) observations, \(p\)
  variables)
\item
  \(\mathbf{X}\) is the design matrix (group indicators)
\item
  \(\mathbf{B}\) is the matrix of coefficients
\item
  \(\mathbf{E}\) is the matrix of errors
\end{itemize}

The null hypothesis is
\(H_{0}:\mathbf{\mu}_{1} = \mathbf{\mu}_{2} = \ldots = \mathbf{\mu}_{g}\)
for all \(p\) variables.

\textbf{ }

\subsection{Test Statistics}

\subsubsection{Wilks' Lambda}

The most commonly used MANOVA test statistic is \textbf{Wilks' Lambda}:

\[\Lambda = |\mathbf{W}\frac{|}{|}\mathbf{T}| = |\mathbf{W}\frac{|}{|}\mathbf{W} + \mathbf{B}|\]

Where:

\begin{itemize}
\item
  \(\mathbf{W}\) is the within-groups sum of squares and cross-products
  matrix
\item
  \(\mathbf{B}\) is the between-groups sum of squares and cross-products
  matrix
\item
  \(\mathbf{T} = \mathbf{W} + \mathbf{B}\) is the total sum of squares
  and cross-products matrix (matrix addition -- note that
  \(|\mathbf{W} + \mathbf{B}| \neq |\mathbf{W}| + |\mathbf{B}|\), since
  determinants are not additive)
\end{itemize}

\textbf{Info: } Wilks' Lambda is a ratio of generalized variances
(determinants), not a literal scalar ``proportion of variance.'' Values
range from 0 to 1; smaller values indicate greater separation between
group mean vectors relative to within-group scatter.

\subsubsection{Other Test Statistics}

Alternative MANOVA test statistics include:

\begin{itemize}
\item
  \textbf{Pillai's Trace}: More robust to violations of assumptions
\item
  \textbf{Hotelling-Lawley Trace}: Most powerful when assumptions are
  met
\item
  \textbf{Roy's Greatest Root}: Tests for difference on the first
  discriminant function only
\end{itemize}

\textbf{ }

\subsection{Assumptions}

For valid MANOVA inference:

\begin{enumerate}
\item
  \textbf{Independence}: Observations must be independent
\item
  \textbf{Multivariate Normality}: Errors follow multivariate normal
  distribution
\item
  \textbf{Homogeneity of Covariance Matrices}:
  \(\mathbf{\Sigma}_{1} = \mathbf{\Sigma}_{2} = \ldots = \mathbf{\Sigma}_{g}\)
\item
  \textbf{Adequate Sample Size}: At least \(n > p + g\) (preferably much
  larger)
\end{enumerate}

\textbf{Tip: } Test homogeneity of covariance matrices using Box's M
test before conducting MANOVA. If violated, consider robust alternatives
or data transformations.

\textbf{ }

\subsection{Interpreting MANOVA Results}

After a significant MANOVA result:

\begin{enumerate}
\item
  \textbf{Examine univariate F-tests} to see which individual variables
  differ across groups
\item
  \textbf{Apply Bonferroni correction} or other multiple comparison
  adjustments
\item
  \textbf{Conduct discriminant analysis} to identify linear combinations
  that best separate groups
\item
  \textbf{Examine effect sizes} (e.g., partial eta-squared) for
  practical significance
\end{enumerate}

\textbf{Example: } Comparing teaching methods (lecture, discussion,
online) on student outcomes measured by exam score, engagement rating,
and satisfaction score. MANOVA tests whether the three methods produce
different patterns across all three outcomes simultaneously.

\textbf{ }

\subsection{Follow-Up Analyses}

After rejecting the null hypothesis in MANOVA:

\begin{itemize}
\item
  \textbf{Univariate ANOVAs} with adjusted alpha levels
\item
  \textbf{Discriminant Function Analysis} to identify separation
  patterns
\item
  \textbf{Post-hoc pairwise comparisons} with appropriate corrections
\item
  \textbf{Contrast analysis} for planned comparisons
\end{itemize}

\textbf{Warning: } Always report which assumptions were tested and any
violations found. MANOVA can be sensitive to assumption violations,
particularly with unequal sample sizes.

\begin{center}\rule{0.5\linewidth}{0.5pt}\end{center}

\textbf{ }

\section{Canonical Correlation Analysis}

\textbf{ }

\subsection{Introduction to Canonical Correlation}

\textbf{Canonical Correlation Analysis (CCA)} examines relationships
between \textbf{two sets of variables} simultaneously. It finds linear
combinations of variables in each set that maximize their correlation.

\textbf{ }

\subsection{Motivation and Applications}

CCA answers questions like:

\begin{itemize}
\item
  How do personality traits (set 1) relate to job performance measures
  (set 2)?
\item
  What is the relationship between physical characteristics and athletic
  performance?
\item
  How do economic indicators relate to social welfare measures?
\end{itemize}

\textbf{Info: } Unlike multiple regression (many predictors, one
response), CCA treats both sets symmetrically and finds the strongest
relationship between linear combinations of each set.

\textbf{ }

\subsection{Mathematical Framework}

Given two sets of variables:

\begin{itemize}
\item
  Set 1: \(\mathbf{X}\) with \(p\) variables
\item
  Set 2: \(\mathbf{Y}\) with \(q\) variables
\end{itemize}

CCA finds coefficients \(\mathbf{a}\) and \(\mathbf{b}\) such that:

\[\text{ Cor}\left( \mathbf{a}'\mathbf{X},\mathbf{b}'\mathbf{Y} \right)\]

is maximized, subject to:

\begin{itemize}
\item
  \(\text{Var}\left( \mathbf{a}'\mathbf{X} \right) = 1\)
\item
  \(\text{Var}\left( \mathbf{b}'\mathbf{Y} \right) = 1\)
\end{itemize}

The resulting linear combinations are called \textbf{canonical
variates}:

\begin{itemize}
\item
  First canonical variate for Set 1:
  \(U_{1} = \mathbf{a}_{1}'\mathbf{X}\)
\item
  First canonical variate for Set 2:
  \(V_{1} = \mathbf{b}_{1}'\mathbf{Y}\)
\end{itemize}

\textbf{ }

\subsection{Number of Canonical Correlations}

The number of canonical correlation pairs equals:

\[k = \min(p,q)\]

Where \(p\) and \(q\) are the number of variables in sets 1 and 2.

Each successive pair of canonical variates:

\begin{itemize}
\item
  Is uncorrelated with previous pairs
\item
  Maximizes correlation subject to this constraint
\item
  Has correlation less than or equal to the previous pair
\end{itemize}

\textbf{Example: } If you have 5 academic performance measures and 3
personality traits, you can extract at most min(5,3) equals 3 canonical
correlation pairs.

\textbf{ }

\subsection{Canonical Loadings}

\textbf{Canonical loadings} (also called \textbf{structure
coefficients}) measure the correlation between original variables and
canonical variates:

\begin{itemize}
\item
  Loading of \(X_{i}\) on \(U_{j}\): \(r_{X_{i},U_{j}}\)
\item
  Loading of \(Y_{i}\) on \(V_{j}\): \(r_{Y_{i},V_{j}}\)
\end{itemize}

These loadings help interpret what each canonical variate represents.

\textbf{Tip: } Canonical loadings are generally more interpretable than
canonical weights (coefficients) because they are not affected by
multicollinearity among predictors.

\textbf{ }

\subsection{Testing Significance}

Test whether there are any significant canonical correlations using
Wilks' Lambda:

\[\Lambda = \prod_{i = 1}^{k}\left( 1 - r_{i}^{2} \right)\]

Where \(r_{i}\) are the canonical correlations.

The test statistic:

\[\chi^{2} = - \left( n - 1 - \frac{p + q + 1}{2} \right)\log(\Lambda)\]

follows approximately a chi-square distribution with \(pq\) degrees of
freedom.

\textbf{ }

\subsection{Interpretation Guidelines}

When interpreting canonical correlations:

\begin{enumerate}
\item
  Test statistical significance of each canonical correlation
\item
  Examine the proportion of variance explained
\item
  Interpret canonical loadings to understand variable contributions
\item
  Name canonical variates based on which variables load highly
\item
  Consider practical significance, not just statistical significance
\end{enumerate}

\textbf{Warning: } Canonical correlation requires large sample sizes for
stable results. A general guideline is at least 10 observations per
variable, preferably 20.

\textbf{ }

\subsection{Relationship to Other Methods}

Canonical correlation generalizes several statistical techniques:

\begin{itemize}
\item
  Multiple regression: Special case where Set 2 has one variable
\item
  Discriminant analysis: When Set 2 is group membership
\end{itemize}

\textbf{Warning: } CCA is \textbf{not} a generalization of PCA. Setting
Set 1 and Set 2 to the same variables makes every canonical correlation
equal to 1 (each variate is trivially perfectly correlated with itself)
and does not recover the principal component directions, which come from
decomposing a single set's own covariance/correlation matrix rather than
the association between two sets.

\textbf{Example: } Examining the relationship between socioeconomic
factors (income, education, housing quality) and health outcomes (life
expectancy, disease prevalence, healthcare access). CCA identifies which
combinations of socioeconomic factors most strongly relate to which
patterns of health outcomes.

\begin{center}\rule{0.5\linewidth}{0.5pt}\end{center}

\textbf{ }

\section{Analysis by Factors and Regression}

\textbf{ }

\subsection{Combining Factor Analysis with Regression}

Factor analysis can be integrated with regression to address
\textbf{multicollinearity} and \textbf{dimensionality reduction} in
predictive modeling.

\textbf{ }

\subsection{The Problem: Multicollinearity}

When predictor variables are highly correlated:

\begin{itemize}
\item
  Regression coefficients become unstable
\item
  Standard errors inflate
\item
  Interpretation becomes difficult
\item
  Prediction may remain good, but inference suffers
\end{itemize}

\textbf{Warning: } High correlation among predictors (typically
\textbar r\textbar{} greater than 0.8) can lead to:

\begin{itemize}
\item
  Nonsensical coefficient estimates
\item
  Coefficients with unexpected signs
\item
  High variance inflation factors (VIF greater than 10)
\end{itemize}

\textbf{ }

\subsection{Factor Scores as Predictors}

The general approach:

\begin{enumerate}
\item
  \textbf{Extract factors} from the predictor variables using factor
  analysis
\item
  \textbf{Compute factor scores} for each observation
\item
  \textbf{Use factor scores} as predictors in regression instead of
  original variables
\end{enumerate}

This addresses multicollinearity because:

\begin{itemize}
\item
  Factor scores are orthogonal (uncorrelated) when using orthogonal
  rotation
\item
  Dimensionality is reduced (fewer predictors than original variables)
\item
  Factors represent underlying constructs rather than individual
  variables
\end{itemize}

\textbf{ }

\subsection{Mathematical Framework}

Given original predictors \(\mathbf{X}\) with \(p\) variables, extract
\(m\) factors (where \(m < p\)):

\[\mathbf{X} = \mathbf{L}\mathbf{F} + \mathbf{U}\]

Where:

\begin{itemize}
\item
  \(\mathbf{L}\) is the loading matrix
\item
  \(\mathbf{F}\) are factor scores
\item
  \(\mathbf{U}\) are unique factors
\end{itemize}

Then fit regression:

\[Y = \beta_{0} + \beta_{1}F_{1} + \beta_{2}F_{2} + \ldots + \beta_{m}F_{m} + \varepsilon\]

\textbf{ }

\subsection{Methods for Computing Factor Scores}

\subsubsection{Regression Method (Thomson)}

For orthogonal factors:

\[\hat{\mathbf{F}} = \mathbf{L}'\mathbf{\Sigma}^{- 1}\mathbf{X}\]

For oblique (correlated) factors with factor correlation matrix
\(\mathbf{\Phi}\):

\[\hat{\mathbf{F}} = \mathbf{\Phi}\mathbf{L}'\mathbf{\Sigma}^{- 1}\mathbf{X}\]

Advantages:

\begin{itemize}
\item
  Maximizes the correlation between true and estimated factor scores
\item
  Most commonly used
\end{itemize}

Note: this is \textbf{not} the same expression as
\(\left( \mathbf{L}'\mathbf{\Sigma}^{- 1}\mathbf{L} \right)^{- 1}\mathbf{L}'\mathbf{\Sigma}^{- 1}\mathbf{X}\),
which is a generalized-least-squares coefficient estimator, not the
regression factor score.

\subsubsection{Bartlett Method}

\[\hat{\mathbf{F}} = \left( \mathbf{L}'\mathbf{\Psi}^{- 1}\mathbf{L} \right)^{- 1}\mathbf{L}'\mathbf{\Psi}^{- 1}\mathbf{X}\]

Where \(\mathbf{\Psi}\) is the diagonal matrix of unique variances.

Advantages:

\begin{itemize}
\item
  Unbiased estimates
\item
  Accounts for measurement error
\end{itemize}

\textbf{Info: } Most statistical software (R, Python, SPSS) provides
options for different factor score estimation methods.

\textbf{ }

\subsection{Benefits of Factor-Based Regression}

\begin{enumerate}
\item
  \textbf{Dimensionality Reduction}: \(m\) factors instead of \(p\)
  variables (where \(m \ll p\))
\item
  \textbf{Elimination of Multicollinearity}: Orthogonal factors have
  zero correlation
\item
  \textbf{Interpretability}: Factors represent meaningful constructs
\item
  \textbf{Robustness}: Less sensitive to measurement error in individual
  variables
\end{enumerate}

\textbf{ }

\subsection{Interpretation Considerations}

When interpreting factor-based regression:

\begin{itemize}
\item
  \textbf{Factor coefficients} indicate importance of underlying
  constructs
\item
  \textbf{Factor loadings} show which original variables contribute to
  each factor
\item
  \textbf{Indirect effects} can be traced from original variables
  through factors to outcome
\end{itemize}

\textbf{Example: } Predicting customer satisfaction from 20 service
quality items. Instead of using all 20 items (with high
multicollinearity), extract 4 factors: Responsiveness, Reliability,
Tangibles, and Empathy. Regress satisfaction on these 4 factors, which
are uncorrelated and more interpretable.

\textbf{ }

\subsection{Limitations and Considerations}

\textbf{Warning: } Factor-based regression has limitations:

\begin{itemize}
\item
  Factors must be meaningful and interpretable
\item
  Factor extraction is somewhat subjective (number of factors, rotation
  method)
\item
  Results depend on the specific sample used for factor analysis
\item
  Prediction in new samples requires computing factor scores using the
  same loadings
\end{itemize}

\textbf{ }

\subsection{Alternative Approaches}

Other methods for addressing multicollinearity:

\begin{itemize}
\item
  \textbf{Ridge Regression}: Shrinks coefficients toward zero
\item
  \textbf{Principal Components Regression}: Similar to factor regression
  but uses PCA
\item
  \textbf{Partial Least Squares}: Finds components that predict Y well
\item
  \textbf{Variable Selection}: Remove redundant predictors
\end{itemize}

\textbf{Tip: } Consider factor-based regression when you have many
correlated predictors representing a smaller number of underlying
constructs. Use PCA-based regression when dimensionality reduction is
the primary goal rather than construct interpretation.

\begin{center}\rule{0.5\linewidth}{0.5pt}\end{center}

\textbf{ }

\section{Programming and Commercial Systems}

\textbf{ }

\subsection{Statistical Software for Multivariate Analysis}

Modern multivariate analysis relies heavily on specialized statistical
software. This section covers the primary tools used in practice and
research.

\textbf{ }

\subsection{Python Ecosystem}

\subsubsection{Core Libraries}

\textbf{NumPy and Pandas}: Foundation for data manipulation

\begin{itemize}
\item
  Matrix operations and linear algebra
\item
  Data structures for multivariate data
\item
  Missing data handling
\end{itemize}

\textbf{Statsmodels}: Statistical modeling

\begin{Shaded}
\begin{Highlighting}[]
\ImportTok{import}\NormalTok{ statsmodels.api }\ImportTok{as}\NormalTok{ sm}
\ImportTok{import}\NormalTok{ statsmodels.formula.api }\ImportTok{as}\NormalTok{ smf}

\CommentTok{\# Logistic regression}
\NormalTok{model }\OperatorTok{=}\NormalTok{ smf.logit(}\StringTok{\textquotesingle{}y \textasciitilde{} x1 + x2 + x3\textquotesingle{}}\NormalTok{, data}\OperatorTok{=}\NormalTok{df)}
\NormalTok{result }\OperatorTok{=}\NormalTok{ model.fit()}

\CommentTok{\# MANOVA}
\ImportTok{from}\NormalTok{ statsmodels.multivariate.manova }\ImportTok{import}\NormalTok{ MANOVA}
\NormalTok{manova }\OperatorTok{=}\NormalTok{ MANOVA.from\_formula(}\StringTok{\textquotesingle{}y1 + y2 + y3 \textasciitilde{} group\textquotesingle{}}\NormalTok{, data}\OperatorTok{=}\NormalTok{df)}
\BuiltInTok{print}\NormalTok{(manova.mv\_test())}
\end{Highlighting}
\end{Shaded}

\textbf{Scikit-learn}: Machine learning and predictive modeling

\begin{Shaded}
\begin{Highlighting}[]
\ImportTok{from}\NormalTok{ sklearn.linear\_model }\ImportTok{import}\NormalTok{ LogisticRegression}
\ImportTok{from}\NormalTok{ sklearn.discriminant\_analysis }\ImportTok{import}\NormalTok{ LinearDiscriminantAnalysis}

\CommentTok{\# Logistic regression with regularization}
\NormalTok{logreg }\OperatorTok{=}\NormalTok{ LogisticRegression(penalty}\OperatorTok{=}\StringTok{\textquotesingle{}l2\textquotesingle{}}\NormalTok{, C}\OperatorTok{=}\FloatTok{1.0}\NormalTok{)}
\NormalTok{logreg.fit(X\_train, y\_train)}

\CommentTok{\# LDA for classification}
\NormalTok{lda }\OperatorTok{=}\NormalTok{ LinearDiscriminantAnalysis()}
\NormalTok{lda.fit(X\_train, y\_train)}
\end{Highlighting}
\end{Shaded}

\subsubsection{Specialized Packages}

\textbf{Factor Analyzer}: Factor analysis

\begin{Shaded}
\begin{Highlighting}[]
\ImportTok{from}\NormalTok{ factor\_analyzer }\ImportTok{import}\NormalTok{ FactorAnalyzer, calculate\_kmo}

\CommentTok{\# KMO test}
\NormalTok{kmo\_all, kmo\_model }\OperatorTok{=}\NormalTok{ calculate\_kmo(df)}

\CommentTok{\# Factor analysis}
\NormalTok{fa }\OperatorTok{=}\NormalTok{ FactorAnalyzer(n\_factors}\OperatorTok{=}\DecValTok{3}\NormalTok{, rotation}\OperatorTok{=}\StringTok{\textquotesingle{}varimax\textquotesingle{}}\NormalTok{)}
\NormalTok{fa.fit(df)}
\NormalTok{loadings }\OperatorTok{=}\NormalTok{ fa.loadings\_}
\end{Highlighting}
\end{Shaded}

\textbf{Scipy}: Scientific computing

\begin{Shaded}
\begin{Highlighting}[]
\ImportTok{from}\NormalTok{ scipy.stats }\ImportTok{import}\NormalTok{ wishart}
\ImportTok{from}\NormalTok{ scipy.spatial.distance }\ImportTok{import}\NormalTok{ mahalanobis}

\CommentTok{\# Wishart distribution sampling}
\NormalTok{W }\OperatorTok{=}\NormalTok{ wishart.rvs(df}\OperatorTok{=}\DecValTok{10}\NormalTok{, scale}\OperatorTok{=}\NormalTok{Sigma, size}\OperatorTok{=}\DecValTok{1}\NormalTok{)}

\CommentTok{\# Mahalanobis distance}
\NormalTok{d }\OperatorTok{=}\NormalTok{ mahalanobis(x, mean, cov\_inv)}
\end{Highlighting}
\end{Shaded}

\textbf{Tip: } Python's ecosystem is ideal for integrating multivariate
analysis with machine learning pipelines and production systems. Use
Jupyter notebooks for exploratory analysis and Python scripts for
production code.

\textbf{ }

\subsection{R Programming}

R provides extensive support for multivariate statistics through base
functions and specialized packages.

\subsubsection{Base R Functions}

\begin{Shaded}
\begin{Highlighting}[]
\CommentTok{\# Hotelling\textquotesingle{}s T{-}squared test}
\FunctionTok{library}\NormalTok{(Hotelling)}
\NormalTok{hotel.test }\OtherTok{\textless{}{-}} \FunctionTok{hotelling.test}\NormalTok{(}\FunctionTok{cbind}\NormalTok{(y1, y2) }\SpecialCharTok{\textasciitilde{}}\NormalTok{ group, }\AttributeTok{data=}\NormalTok{df)}

\CommentTok{\# MANOVA}
\NormalTok{manova\_result }\OtherTok{\textless{}{-}} \FunctionTok{manova}\NormalTok{(}\FunctionTok{cbind}\NormalTok{(y1, y2, y3) }\SpecialCharTok{\textasciitilde{}}\NormalTok{ group, }\AttributeTok{data=}\NormalTok{df)}
\FunctionTok{summary}\NormalTok{(manova\_result, }\AttributeTok{test=}\StringTok{"Wilks"}\NormalTok{)}

\CommentTok{\# Canonical correlation}
\FunctionTok{library}\NormalTok{(CCA)}
\NormalTok{cc\_result }\OtherTok{\textless{}{-}} \FunctionTok{cc}\NormalTok{(X, Y)}
\end{Highlighting}
\end{Shaded}

\subsubsection{Specialized Packages}

\begin{itemize}
\item
  \texttt{mvtnorm}: Multivariate normal and t distributions
\item
  \texttt{MVN}: Multivariate normality tests
\item
  \texttt{car}: Companion to Applied Regression (includes MANOVA)
\item
  \texttt{FactoMineR}: Multivariate exploratory analysis
\item
  \texttt{psych}: Psychological methods including factor analysis
\end{itemize}

\textbf{Info: } R excels at statistical analysis and has more
specialized multivariate packages than Python. It is widely used in
academic research and has excellent documentation.

\textbf{ }

\subsection{Commercial Software Packages}

\subsubsection{SPSS (IBM)}

\textbf{Graphical User Interface}:

\begin{itemize}
\item
  Point-and-click menus for common analyses
\item
  Integrated data editor and output viewer
\item
  Extensive built-in help and tutorials
\end{itemize}

\textbf{Syntax Programming}:

\begin{Shaded}
\begin{Highlighting}[]
\NormalTok{* Logistic regression}
\NormalTok{LOGISTIC REGRESSION VAR=outcome}
\NormalTok{  /METHOD=ENTER predictor1 predictor2 predictor3}
\NormalTok{  /PRINT=GOODFIT CI(95).}

\NormalTok{* MANOVA}
\NormalTok{MANOVA y1 y2 y3 BY group(1,3)}
\NormalTok{  /PRINT=CELLINFO(MEANS)}
\NormalTok{  /DESIGN.}

\NormalTok{* Factor analysis}
\NormalTok{FACTOR /VARIABLES var1 TO var10}
\NormalTok{  /EXTRACTION PC}
\NormalTok{  /ROTATION VARIMAX.}
\end{Highlighting}
\end{Shaded}

Strengths:

\begin{itemize}
\item
  User-friendly for non-programmers
\item
  Comprehensive documentation
\item
  Wide adoption in social sciences
\end{itemize}

\subsubsection{SAS (Statistical Analysis System)}

\textbf{Procedures for Multivariate Analysis}:

\begin{Shaded}
\begin{Highlighting}[]
\NormalTok{/* Logistic regression */}
\NormalTok{PROC LOGISTIC DATA=mydata;}
\NormalTok{  MODEL outcome(EVENT=\textquotesingle{}1\textquotesingle{}) = x1 x2 x3 / LINK=LOGIT;}
\NormalTok{  OUTPUT OUT=pred PREDICTED=prob;}
\NormalTok{RUN;}

\NormalTok{/* MANOVA */}
\NormalTok{PROC GLM DATA=mydata;}
\NormalTok{  CLASS group;}
\NormalTok{  MODEL y1 y2 y3 = group;}
\NormalTok{  MANOVA H=group / PRINTH PRINTE;}
\NormalTok{RUN;}

\NormalTok{/* Canonical correlation */}
\NormalTok{PROC CANCORR DATA=mydata;}
\NormalTok{  VAR x1 x2 x3;}
\NormalTok{  WITH y1 y2 y3;}
\NormalTok{RUN;}

\NormalTok{/* Factor analysis */}
\NormalTok{PROC FACTOR DATA=mydata METHOD=PRINCIPAL}
\NormalTok{  ROTATE=VARIMAX NFACTORS=3;}
\NormalTok{  VAR var1{-}var10;}
\NormalTok{RUN;}
\end{Highlighting}
\end{Shaded}

Strengths:

\begin{itemize}
\item
  Enterprise-level data handling
\item
  Excellent performance with large datasets
\item
  Strong in clinical trials and regulatory environments
\end{itemize}

\subsubsection{MATLAB}

MATLAB's Statistics and Machine Learning Toolbox includes:

\begin{itemize}
\item
  \texttt{fitglm}: Generalized linear models (logistic regression)
\item
  \texttt{manova1}: One-way MANOVA
\item
  \texttt{canoncorr}: Canonical correlation analysis
\item
  \texttt{factoran}: Factor analysis
\end{itemize}

Strengths:

\begin{itemize}
\item
  Matrix operations are native
\item
  Integration with engineering applications
\item
  Excellent for simulation studies
\end{itemize}

\textbf{Warning: } Commercial software (SPSS, SAS, MATLAB) requires
licenses that can be expensive. Academic institutions often provide
access. Consider open-source alternatives (R, Python) for cost-effective
solutions with comparable functionality.

\textbf{ }

\subsection{Software Selection Guidelines}

Choose software based on:

\begin{enumerate}
\item
  \textbf{Task Requirements}:

  \begin{itemize}
  \item
    Exploratory analysis: R or SPSS
  \item
    Production systems: Python
  \item
    Large-scale data: SAS
  \item
    Engineering integration: MATLAB
  \end{itemize}
\item
  \textbf{Team Expertise}:

  \begin{itemize}
  \item
    Programmers: Python, R
  \item
    Non-programmers: SPSS
  \item
    Mixed teams: R + RStudio
  \end{itemize}
\item
  \textbf{Budget Constraints}:

  \begin{itemize}
  \item
    Free: R, Python
  \item
    Institutional license: SPSS, SAS, MATLAB
  \item
    Cloud-based: Python (widely available on cloud platforms)
  \end{itemize}
\item
  \textbf{Reproducibility Needs}:

  \begin{itemize}
  \item
    Code-based: Python, R, SAS, MATLAB
  \item
    Documentation: All have good capabilities
  \end{itemize}
\end{enumerate}

\textbf{Example: } A pharmaceutical company might use:

\begin{itemize}
\item
  SAS for regulatory submissions (required by FDA)
\item
  R for exploratory analysis and visualization
\item
  Python for machine learning and predictive modeling
\item
  SPSS for surveys and marketing research
\end{itemize}

\textbf{ }

\subsection{Best Practices}

Regardless of software choice:

\begin{enumerate}
\item
  \textbf{Document your analysis}: Use scripts/syntax, not just
  point-and-click
\item
  \textbf{Version control}: Track changes to code and data
\item
  \textbf{Reproducibility}: Include random seeds, package versions
\item
  \textbf{Validation}: Cross-check critical results across software when
  possible
\item
  \textbf{Backup}: Maintain copies of data, code, and results
\end{enumerate}

\textbf{Tip: } Learn one statistical language well (Python or R) and
become familiar with at least one commercial package. This combination
provides flexibility for different work environments and requirements.

\begin{center}\rule{0.5\linewidth}{0.5pt}\end{center}

\textbf{ }

\section{Practice Questions and Answers}

\textbf{ }

\subsection{Question 1: Logistic Regression}

\textbf{Question:} Why is the logit link function necessary in logistic
regression?

\textbf{Answer:} The logit link function \(\log(\frac{p}{1 - p})\)
transforms probabilities (bounded between 0 and 1) to the real line
(negative infinity to positive infinity). This allows us to model
probabilities using a linear combination of predictors. Without this
transformation, linear regression would produce predicted probabilities
outside the valid range of 0,1, leading to nonsensical results. The
inverse logit (logistic function) then maps predictions back to valid
probability scale.

\textbf{ }

\subsection{Question 2: Wishart Distribution}

\textbf{Question:} What is the relationship between the Wishart
distribution and the chi-square distribution?

\textbf{Answer:} The Wishart distribution is the multivariate
generalization of the chi-square distribution. Just as the chi-square
distribution describes the distribution of a sum of squared standard
normal variables, the Wishart distribution describes the distribution of
sample covariance matrices from multivariate normal data. When \(p = 1\)
(univariate case), the Wishart distribution reduces to the chi-square
distribution. This relationship is fundamental for conducting inference
about covariance matrices in the multivariate setting.

\textbf{ }

\subsection{Question 3: Hotelling's T-Squared}

\textbf{Question:} When would you use Hotelling's \(T^{2}\) test instead
of multiple univariate t-tests?

\textbf{Answer:} Use Hotelling's \(T^{2}\) when comparing mean vectors
across groups on multiple correlated variables simultaneously.
Conducting separate t-tests inflates Type I error (the family-wise error
rate increases with each test). Hotelling's \(T^{2}\) provides a single
omnibus test that controls overall Type I error while accounting for
correlations between variables. After rejecting with Hotelling's
\(T^{2}\), you can conduct follow-up univariate tests with appropriate
corrections to identify which specific variables differ.

\textbf{ }

\subsection{Question 4: MANOVA vs Multiple ANOVAs}

\textbf{Question:} Explain why MANOVA is preferred over conducting
separate ANOVAs on each dependent variable.

\textbf{Answer:} MANOVA is preferred for three main reasons:

\begin{enumerate}
\item
  \textbf{Type I Error Control}: Multiple ANOVAs increase the
  probability of false positives through multiple testing. MANOVA
  maintains the overall Type I error rate at the specified alpha level.
\item
  \textbf{Accounting for Correlations}: MANOVA incorporates the
  correlation structure among dependent variables, potentially
  increasing statistical power.
\item
  \textbf{Detecting Patterns}: MANOVA can detect differences in linear
  combinations of variables that might be missed by individual ANOVAs,
  revealing multivariate patterns of group differences.
\end{enumerate}

\textbf{ }

\subsection{Question 5: Wilks' Lambda Interpretation}

\textbf{Question:} What does Wilks' Lambda measure and how is it
interpreted in MANOVA?

\textbf{Answer:} Wilks' Lambda measures the proportion of total variance
in the dependent variables that is not explained by differences between
groups. It is calculated as \(\Lambda = |W\frac{|}{|}T|\) where \(W\) is
the within-groups matrix and \(T\) is the total matrix. Values range
from 0 to 1:

\begin{itemize}
\item
  \(\Lambda\) near 0 indicates strong group separation (most variance is
  between-groups)
\item
  \(\Lambda\) near 1 indicates weak group separation (most variance is
  within-groups)
\end{itemize}

Smaller values of Wilks' Lambda lead to rejection of the null hypothesis
that group means are equal.

\textbf{ }

\subsection{Question 6: Canonical Correlation}

\textbf{Question:} How many canonical correlation pairs can be extracted
when analyzing the relationship between 7 socioeconomic variables and 4
health outcomes?

\textbf{Answer:} The number of canonical correlation pairs equals
\(\min(p,q)\) where \(p\) and \(q\) are the number of variables in each
set. In this case, \(\min(7,4) = 4\), so exactly 4 canonical correlation
pairs can be extracted. The first pair has the maximum possible
correlation, and each successive pair has decreasing correlation while
being uncorrelated with previous pairs.

\textbf{ }

\subsection{Question 7: Canonical Loadings}

\textbf{Question:} What are canonical loadings and why are they
important for interpretation?

\textbf{Answer:} Canonical loadings (structure coefficients) represent
the correlation between original variables and canonical variates. For
example, the loading of variable \(X_{3}\) on canonical variate
\(U_{1}\) measures how strongly \(X_{3}\) contributes to \(U_{1}\).
These loadings are crucial for interpretation because they help identify
what each canonical variate represents. Variables with high loadings
(typically \textbar r\textbar{} greater than 0.3 or 0.4) on a variate
are considered important contributors to that dimension. Canonical
loadings are generally more interpretable than canonical weights because
they are not affected by multicollinearity.

\textbf{ }

\subsection{Question 8: Factor Scores in Regression}

\textbf{Question:} What are the main benefits of using factor scores as
predictors in regression analysis?

\textbf{Answer:} Using factor scores as predictors provides several
benefits:

\begin{enumerate}
\item
  \textbf{Multicollinearity Reduction}: Factor scores are orthogonal
  (uncorrelated) when using orthogonal rotation, eliminating
  multicollinearity problems
\item
  \textbf{Dimensionality Reduction}: Fewer predictors than original
  variables (\(m\) factors instead of \(p\) variables where \(m \ll p\))
\item
  \textbf{Construct Interpretation}: Factors represent underlying latent
  constructs rather than individual observed variables
\item
  \textbf{Robustness}: Less sensitive to measurement error in individual
  variables since factors capture shared variance
\end{enumerate}

This approach is particularly useful when you have many correlated
predictors representing a smaller number of theoretical constructs.

\textbf{ }

\subsection{Question 9: Box's M Test}

\textbf{Question:} What does Box's M test assess and why is it important
for MANOVA?

\textbf{Answer:} Box's M test assesses the assumption of homogeneity of
covariance matrices across groups, testing
\(H_{0}:\Sigma_{1} = \Sigma_{2} = \ldots = \Sigma_{g}\). This is crucial
for MANOVA because the method assumes groups have equal covariance
structures. Violation of this assumption can affect the validity of
MANOVA results, particularly with unequal sample sizes. However, Box's M
is sensitive to departures from multivariate normality, so it should be
used cautiously. If the test is significant, consider robust
alternatives or data transformations.

\textbf{ }

\subsection{Question 10: Software Selection}

\textbf{Question:} What factors should guide the choice between Python's
statsmodels and SPSS for conducting multivariate regression analysis in
a corporate setting?

\textbf{Answer:} Key factors include:

\begin{enumerate}
\item
  \textbf{Team Skills}: SPSS is more suitable for teams without
  programming background (GUI-based), while Python requires coding
  skills
\item
  \textbf{Integration Needs}: Python integrates easily with databases,
  web applications, and machine learning pipelines; SPSS is more
  standalone
\item
  \textbf{Cost}: Python is free and open-source; SPSS requires expensive
  licenses
\item
  \textbf{Reproducibility}: Both support scripting (Python natively,
  SPSS through syntax), crucial for reproducible research
\item
  \textbf{Advanced Features}: Python offers more flexibility for custom
  analyses and newer methods; SPSS has more built-in standard procedures
\end{enumerate}

In practice, many organizations use both: SPSS for standard analyses by
non-programmers, Python for custom work and production systems.

\begin{center}\rule{0.5\linewidth}{0.5pt}\end{center}

\textbf{ }

\section{Summary and Key Takeaways}

\textbf{ }

\subsection{Logistic Regression}

\begin{itemize}
\item
  Models binary outcomes using Bernoulli distribution
\item
  Logit link function: \(\log(\frac{p}{1 - p})\)
\item
  Maximum likelihood estimation (not OLS)
\item
  Coefficients represent change in log-odds
\end{itemize}

\textbf{ }

\subsection{Covariance Matrix Inference}

\begin{itemize}
\item
  Based on multivariate normal distribution
\item
  Wishart distribution generalizes chi-square
\item
  Box's M test for equality of covariances
\item
  Bartlett's test for univariate variances
\end{itemize}

\textbf{ }

\subsection{Mean Vector Inference}

\begin{itemize}
\item
  Hotelling's \(T^{2}\) generalizes t-test
\item
  Can be transformed to F distribution
\item
  Controls Type I error for multiple variables
\item
  Applicable to one-sample, two-sample, and paired designs
\end{itemize}

\textbf{ }

\subsection{MANOVA}

\begin{itemize}
\item
  Tests multiple dependent variables simultaneously
\item
  Controls Type I error inflation
\item
  Wilks' Lambda most common test statistic
\item
  Requires multivariate normality and homogeneous covariances
\end{itemize}

\textbf{ }

\subsection{Canonical Correlation}

\begin{itemize}
\item
  Examines relationships between two variable sets
\item
  Maximizes correlation between linear combinations
\item
  Number of pairs equals \(\min(p,q)\)
\item
  Canonical loadings aid interpretation
\end{itemize}

\textbf{ }

\subsection{Factor-Based Regression}

\begin{itemize}
\item
  Addresses multicollinearity through dimensionality reduction
\item
  Factor scores used as predictors
\item
  Orthogonal factors have zero correlation
\item
  Balances interpretability and parsimony
\end{itemize}

\textbf{ }

\subsection{Statistical Software}

\begin{itemize}
\item
  \textbf{Python}: statsmodels, scikit-learn, flexible integration
\item
  \textbf{R}: Comprehensive multivariate packages, research-oriented
\item
  \textbf{SPSS}: User-friendly GUI, social science standard
\item
  \textbf{SAS}: Enterprise data handling, regulatory compliance
\item
  Choose based on team skills, requirements, and budget
\end{itemize}

\begin{center}\rule{0.5\linewidth}{0.5pt}\end{center}

\textbf{ }

\section{References and Further Reading}

\textbf{ }

\subsection{Recommended Textbooks}

\begin{enumerate}
\item
  Johnson, R. A., \& Wichern, D. W. (2007). \textbf{Applied Multivariate
  Statistical Analysis} (6th ed.). Pearson.
\item
  Rencher, A. C., \& Christensen, W. F. (2012). \textbf{Methods of
  Multivariate Analysis} (3rd ed.). Wiley.
\item
  Hair, J. F., Black, W. C., Babin, B. J., \& Anderson, R. E. (2018).
  \textbf{Multivariate Data Analysis} (8th ed.). Cengage.
\item
  Tabachnick, B. G., \& Fidell, L. S. (2018). \textbf{Using Multivariate
  Statistics} (7th ed.). Pearson.
\end{enumerate}

\textbf{ }

\subsection{Online Resources}

\begin{itemize}
\item
  Python statsmodels documentation: \url{https://www.statsmodels.org/}
\item
  Scikit-learn user guide:
  \url{https://scikit-learn.org/stable/user_guide.html}
\item
  R CRAN Task View Multivariate:
  \url{https://cran.r-project.org/web/views/Multivariate.html}
\item
  Penn State STAT 505 (Online course on multivariate methods)
\end{itemize}

\textbf{ }

\subsection{Software Documentation}

\begin{itemize}
\item
  \textbf{Python}: pandas, numpy, scipy, statsmodels, scikit-learn
\item
  \textbf{R}: Base R stats package, car, MVN, FactoMineR, psych
\item
  \textbf{SPSS}: IBM SPSS Statistics documentation
\item
  \textbf{SAS}: SAS/STAT User's Guide
\end{itemize}

\textbf{End of Lecture Notes}

MA2003B - Multivariate Regression

Tec de Monterrey

\end{document}

